\documentclass[a4,10pt]{article}

\usepackage{xeCJK}
\usepackage{mathptmx}
\usepackage{anyfontsize}
\usepackage{t1enc}
\usepackage{graphicx}
\usepackage[top=20mm,bottom=20mm,left=20mm,right=20mm]{geometry}
\usepackage{wrapfig}
\usepackage{float} %设置图片浮动位置的宏包
\usepackage{subfigure} %插入多图时用子图显示的宏包
\usepackage[justification=centering]{caption}
\usepackage{ctex}
\usepackage[
backend=biber,
style=nature,
]{biblatex}

\addbibresource{citation.bib}

\makeatletter
  \def\vhrulefill#1{\leavevmode\leaders\hrule\@height#1\hfill \kern\z@}
\makeatother

\pagestyle{plain}

\begin{document}
\begin{titlepage}
    \begin{center}
        \includegraphics[scale=0.5]{Figure/SCU.png}\\
        

        \vspace{1.5cm}

        \textbf{\zihao{1}\songti{本\hspace{1cm}科\hspace{1cm}实\hspace{1cm}验\hspace{1cm}报\hspace{1cm}告}}\\

        \vspace{1cm}

        \includegraphics[scale=0.5]{Figure/SCU_logo.png}\\

        \vspace{2cm}

        \begin{tabular}{ll}
            \zihao{4}\songti{课\hspace{0.5em}程\hspace{0.5em}名\hspace{0.5em}称：\underline{动物生物学实验\hspace{1em}}} & \zihao{4}\songti{课\hspace{1.25em}程\hspace{1.25em}号：\underline{2041xxxxx\hspace{3.5em}}} \\
            \\
            \zihao{4}\songti{所\hspace{0.5em}在\hspace{0.5em}学\hspace{0.5em}院：\underline{生命科学院\hspace{3em}}} & \zihao{4}\songti{专\hspace{3.5em}业：\underline{生物科学\hspace{4em}}} \\
            \\
            \zihao{4}\songti{学\hspace{0.5em}生\hspace{0.5em}姓\hspace{0.5em}名：\underline{何慕菱\hspace{5em}}} & \zihao{4}\songti{学\hspace{3.5em}号：\underline{2021xxxxxxxxx\hspace{1.5em}}}\\
            \\
            \zihao{4}\songti{课\hspace{1.25em}序\hspace{1.25em}号：\underline{02\hspace{7em}}} & \zihao{4}\songti{年\hspace{3.5em}级：\underline{2022级\hspace{5em}}}\\
            \\
            \zihao{4}\songti{指\hspace{0.5em}导\hspace{0.5em}教\hspace{0.5em}师：\underline{XX\hspace{6em}}} & \zihao{4}\songti{成\hspace{0.5em}绩\hspace{0.5em}评\hspace{0.5em}定：\underline{\hspace{8em}}}
            \\
        \end{tabular}

        \vspace{5.8cm}

        \zihao{4}\songti{二〇二三年五月}

    \end{center}
\end{titlepage}

% 标题区域开始
\begin{center}
    \zihao{2}四川大学实验报告 \\
    \vspace{1em}
    \begin{tabular}{lll}
        \zihao{4}学\hspace{1em}院\hspace{1em}生命科学院 & \zihao{4}专\hspace{2em}业\hspace{1em}生物科学\hspace{2em} & \zihao{4}2022级管理0班\hspace{1em}组\\
  
        \zihao{4}姓\hspace{1em}名\hspace{1em}何慕菱 & \zihao{4}同实验者\hspace{1em}小红 & \zihao{4}22年10月24日
    \end{tabular}
    
    \vhrulefill{2pt}
\end{center}
\zihao{4}题\hspace{1em}目：XX的外形观察与内部解刨
% 标题区域结束

\section{实验目的}
\large{} 
\begin{enumerate}
  \item 学习XX的解刨方法
  \item 通过解刨认识、学习XX的内部结构
\end{enumerate}

\section{实验原理}
这里写实验原理

\section{实验步骤}
\begin{enumerate}
  \item 第一步
  \item 第二部
  \item etc...
\end{enumerate}

\section{实验结果}
描述实验现象

\section{讨论}

\subsection{XX与温度的关系}

\textbf{实验方法：}

\textbf{数据记录：}
\begin{table}[h]
  \centering
  \begin{tabular}{|c|c|c|c|c|c|c|c|}
  \hline
  温度$(^\circ \mathrm{C})$ & 室温 & \makebox[2.5em]{4} & 5$\sim$16 & 16$\sim$20 & 20$\sim$36 & 36$\sim$40 & 40$\sim$47 \\ \hline
  心率(次/分) & 6  & 2 & 3    & 5     & 8     & 2     & 0     \\ \hline
  \end{tabular}
\end{table}

\textbf{现象记录：}
描述现象

\textbf{心率-温度曲线：}
\vspace{1em}

\begin{minipage}{16em}
  \includegraphics[scale=0.38]{Figure/output_plot.png}
\end{minipage}
\begin{minipage}{23em}
  描述图标内容
\end{minipage}

\subsection{血液循环途径的研究方法}
对相关内容进一步讨论。

同时在讨论中适当引用\cite{badea2008vivo,bigby1988harrison}。

XXX...\cite{loscalzo2022harrison}

因此，...

\clearpage
\appendix
\section{图片}

在这里插入图片



\clearpage
\printbibliography

\end{document}